% AI-written, AI cross-reviewed; no human sign-off; not a claim that RH is proved
% Atlas unit T22. Local manuscript v3, item 5.32, heading on printed PDF p. 47.
% UpperTheory.LiteratureTheorems.zarhin_statement
% paper2/split_p012.tex:387–399. Verbatim excerpt; SHA-256 2adab53bac70e3ca610fa73c36b99e1da508bf6d2d477619a58145bf11e7ce63.
% Original TeX uses the macros and cross-reference keys of the source manuscript.


% Start of atlas edition annotation. The following block is outside the original excerpt.
\paragraph{Status of named inputs in this reading.}
\noindent\nolinkurl{PositivityTheory.HtPt.Narrow.zarhin} --- in-repo constant (\path{Development/WeilBound/CoreStatements/HodgeTatePoints/ZarTop.lean:6--7}); not compared.\par
\noindent\nolinkurl{UpperTheory.LiteratureTheorems.zarhin_statement} --- in-repo constant (\path{Development/AutomorphicArithmetic/ThetaCorrespondence/ThetaPairLifts/ThirdPlaceSeries/EveryCharacterOfLiteratureOfVWin.lean:547--548}); not compared.\par
\noindent\nolinkurl{PositivityTheory.PrwZar.zarhinStatement_of_dual_cube} --- in-repo constant (\path{Development/WeilBound/CoreStatements/AbelianHeightsDuality/ZarhinTrick/ZarhinStatementOfDualCubeQiR.lean:68--70}); not compared.\par
\noindent\nolinkurl{PositivityTheory.PrwZar.zarhinStatement_of_dual_cube_qiR} --- in-repo constant (\path{Development/WeilBound/CoreStatements/AbelianHeightsDuality/ZarhinTrick/ZarhinStatementOfDualCubeQiR.lean:45--66}); not compared.\par
% End of atlas edition annotation.

\begin{formalstatement}[Splitting off an abelian scheme]\label{sub:T22}

Given the open-literature record, let $C$ be algebraically closed with a
nonnegative-real-valued valuation and let $R\to C$ identify $R$ as its
valuation ring. Every abelian scheme $A/R$ with a relatively ample
invertible sheaf admits group homomorphisms $i:A\to P$, $r:P\to A$
with $r i=1_A$. The scheme $P/R$ is abelian, smooth of a specified
relative dimension and carries a rigidified symmetric relatively ample
invertible sheaf with trivial polarization kernel.

\end{formalstatement}\subnode{T22}

\noindent The cited theorem works over a field. The formal target uses a valuation-ring base and a symmetric rigidified relatively ample line. \href{https://www.jmilne.org/math/CourseNotes/AVc.pdf}{J.S. Milne, Abelian Varieties, Theorem 13.12}~\cite{r21-milne-av}.
